\section{技术创新}
\subsection{核心技术创新点}

本方案围绕嵌入式推理开发与部署，拟研究统一算子语义、可核验的三后端执行，以及静态执行计划与跨操作系统移植。\ProcessorCore{}、\VectorUnit{}和\AcceleratorName{}提供计算能力；项目拟增加模型定义、后端选择与验证记录之间的统一衔接。这是\TechnicalBaseline{}提出的研究方向，改进效果需通过对照实验检验。\cite{technicalreport}

\textbf{统一算子语义，约束转换与执行的一致性。}算子描述记录数学运算、形状、布局、步长、数据类型和量化参数。首个算例限定对称有符号INT8、每张量量化、零点为0，明确舍入、饱和及累加范围。受限转换器在电脑端检查约束，生成算子顺序、权重、缓冲区计划和后端选择；不支持的算子在转换阶段报告。这为接口衔接和逐层结果核对提供共同依据。

\textbf{以独立参考验证标量、向量与矩阵后端。}CPU标量实现提供功能基线，RVV处理适合向量化的计算，NPU承担支持的规整矩阵运算。对同一算子的可用后端，以相同输入和量化定义逐元素比较。独立参考采用INT64累加并检查范围；NPU回读完整INT32矩阵结果，再完成偏置、激活与重定标。用例覆盖尾部元素、饱和边界和舍入规则，避免最终分类相同掩盖中间计算错误。

\textbf{用静态执行计划连接开发与系统移植。}\RuntimeName{}按预定顺序调度算子与工作缓冲区，通过普通C指针和参数结构连接后端，RVV类型限制在后端内部。首期NPU采用单所有者串行提交与完成等待，分别检查搬运、翻译状态和一致性。推理核心源码在\RealTimeSystem{}与Linux中复用，应用及镜像分别构建、重启切换；以结果一致性、适配代码范围和重复运行记录评价移植效果。

实验分别记录转换检查、逐层误差、端到端正确性、运行时间和存储需求。量化模型相对浮点模型的识别效果单独评估，不将整数后端一致性等同于模型准确率。

\subsection{自主知识产权}

\AcceleratorName{}已有硬件生成与C接口，Merlin等已有算子生成与映射研究，均属于本项目借鉴和复用的开源基础。使用这些计算单元不构成原创声明，上游映射方法不能表述为本团队首次提出。\cite{gemmini,merlin}

拟形成的自主工作包括受限模型转换、算子注册与参数检查、独立参考实现、后端适配、静态运行时和测试工具。按实际实现记录作者、版本与修改内容，维护组件来源、许可证及使用范围，区分团队贡献和上游能力。当前不列尚未取得的专利、软件著作权或测试证明。
